\chapter{Linear Models, GLMs and SVM Basics}\label{n:linear}
\studytime{2.5 hours}

\section{Linear Regression}
\begin{dfn}{Linear regression}
\textbf{Linear regression} is a supervised learning method that models a numerical target as a \textbf{linear combination} of the input features plus an intercept:
$\ \hat y=\theta_0+\theta_1x_1+\dots+\theta_dx_d=\vect\theta^{\!\top}\vect x$.
The parameters are chosen to minimise the \textbf{sum of squared errors} (least squares).
\end{dfn}

\textbf{Cost function:} $J(\vect\theta)=\frac1{2n}\sum_{i=1}^n(\hat y_i-y_i)^2$ (MSE up to a factor).

\textbf{Two ways to find $\vect\theta$:}
\begin{itemize}
  \item \textbf{Closed form (normal equation):} $\vect\theta=(\mathbf X^{\!\top}\mathbf X)^{-1}\mathbf X^{\!\top}\vect y$. For one feature: $\theta_1=S_{xy}/S_{xx}$, $\theta_0=\bar y-\theta_1\bar x$. Exact; slow when $d$ is large ($O(d^3)$).
  \item \textbf{Gradient descent:} iterative; works for huge data.
\end{itemize}

\begin{numer}[: least-squares line and $R^2$]
Hours $x=(1,2,3,4,5)$, marks $y=(35,45,50,65,70)$; $\bar x=3$, $\bar y=53$.\\
$S_{xy}=36+8+0+12+34=90$, $S_{xx}=10$ $\Rightarrow$ $\theta_1=9$, $\theta_0=53-27=26$. \textbf{Model:} $\hat y=26+9x$; $x=6\Rightarrow\hat y=80$.\\
$\hat y=(35,44,53,62,71)$, residuals $(0,1,-3,3,-1)$, SSE $=20$, SST $=\sum(y-\bar y)^2=830$.\\
MSE $=4$, RMSE $=2$, $R^2=1-\frac{20}{830}=\mathbf{0.976}$ (the line explains 97.6\% of the variation).
\end{numer}

\begin{dfn}{Coefficient of determination $R^2$}
$R^2=1-\dfrac{\mathrm{SSE}}{\mathrm{SST}}$: the fraction of the variance of $y$ explained by the model ($0\le R^2\le1$ on training data). \textbf{Adjusted} $R^2=1-(1-R^2)\frac{n-1}{n-p-1}$ penalises extra features.
\end{dfn}

\subsection{Gradient descent}
\begin{dfn}{Gradient descent}
\textbf{Gradient descent} is an iterative optimisation algorithm that minimises a cost function by repeatedly moving the parameters a small step in the direction of the \textbf{negative gradient}:
$\ \theta_j\leftarrow\theta_j-\alpha\,\dfrac{\partial J}{\partial\theta_j}$, where $\alpha$ is the \textbf{learning rate}.
\end{dfn}
For linear regression: $\frac{\partial J}{\partial\theta_0}=\frac1n\sum(\hat y_i-y_i)$, $\ \frac{\partial J}{\partial\theta_1}=\frac1n\sum(\hat y_i-y_i)x_i$.
\begin{itemize}
  \item $\alpha$ too small: very slow; $\alpha$ too large: overshoots and diverges.
  \item \textbf{Batch GD} uses all data per step (smooth, slow); \textbf{Stochastic GD} one example per step (fast, noisy, escapes local minima); \textbf{Mini-batch GD} small batches (best of both, used on GPUs).
  \item Scale features first, otherwise the cost bowl is elongated and GD is slow.
\end{itemize}

\begin{numer}[: one GD step]
Marks data, start $\theta_0=\theta_1=0$, $\alpha=0.01$. All $\hat y=0$.\\
$\frac{\partial J}{\partial\theta_0}=\frac{-(35+45+50+65+70)}{5}=-53$;\ $\frac{\partial J}{\partial\theta_1}=\frac{-(35+90+150+260+350)}{5}=-177$.\\
$\theta_0=0-0.01(-53)=0.53$,\ $\theta_1=0-0.01(-177)=1.77$.
\end{numer}

\begin{drawit}{Gradient descent on a cost curve}
\centering
\begin{tikzpicture}
\begin{axis}[width=7cm,height=4.2cm,xlabel={$\theta$},ylabel={$J(\theta)$},ticks=none,xmin=-3.2,xmax=3.2,ymin=0,ymax=10]
\addplot[thick,mlblue,domain=-3:3]{x^2};
\addplot[mlred,mark=*,thick,-Stealth] coordinates {(-2.8,7.84)(-1.9,3.61)(-1.2,1.44)(-0.7,0.49)(-0.35,0.12)};
\node[font=\scriptsize] at (axis cs:1.5,8.5) {steps get smaller near minimum};
\end{axis}
\end{tikzpicture}
\end{drawit}

\subsection{Polynomial regression and regularisation}
\begin{itemize}
  \item \textbf{Polynomial regression:} add powers of features ($x^2,x^3,\dots$) and fit a linear model; high degree $\Rightarrow$ overfitting.
  \item \textbf{Ridge (L2):} $J+\lambda\sum\theta_j^2$, shrinks weights.
  \item \textbf{Lasso (L1):} $J+\lambda\sum|\theta_j|$, sets some weights to zero (feature selection).
  \item \textbf{Elastic Net:} mix of L1 and L2. Larger $\lambda$ $\Rightarrow$ simpler model.
\end{itemize}

\apa{simple, fast, interpretable coefficients, closed-form solution}{assumes linearity, sensitive to outliers and multicollinearity}{price prediction, sales forecasting, trend analysis, risk estimation}

\section{Logistic Regression}
\begin{dfn}{Logistic regression}
\textbf{Logistic regression} is a supervised \textbf{classification} algorithm that estimates the probability that an instance belongs to the positive class by passing a linear combination of the features through the \textbf{sigmoid} function:
$\ p=\sigma(\vect\theta^{\!\top}\vect x)=\dfrac1{1+e^{-\vect\theta^{\!\top}\vect x}}$. Predict 1 if $p\ge0.5$, else 0.
\end{dfn}

\begin{itemize}
  \item Sigmoid maps any $z$ into $(0,1)$; $\sigma(0)=0.5$; $\sigma'(z)=\sigma(z)(1-\sigma(z))$.
  \item \textbf{Odds} $=\frac p{1-p}$; \textbf{log-odds (logit)} $=\ln\frac p{1-p}=\vect\theta^{\!\top}\vect x$ (linear in $x$). One unit increase in $x_j$ multiplies the odds by $e^{\theta_j}$.
  \item \textbf{Decision boundary:} $\vect\theta^{\!\top}\vect x=0$ (a straight line/plane).
  \item \textbf{Cost (log-loss / binary cross-entropy):} $J=-\frac1n\sum\big[y_i\ln p_i+(1-y_i)\ln(1-p_i)\big]$. Convex; no closed form; minimise with GD. Gradient $\frac1n\sum(p_i-y_i)x_{ij}$.
  \item \textbf{Softmax regression} (multinomial): $P(y=k)=\frac{e^{z_k}}{\sum_je^{z_j}}$ for more than two classes.
\end{itemize}

\begin{drawit}{Sigmoid curve}
\centering
\begin{tikzpicture}
\begin{axis}[width=7cm,height=4cm,xlabel={$z$},ylabel={$\sigma(z)$},xmin=-6,xmax=6,ymin=0,ymax=1,ytick={0,0.5,1},xtick={-6,0,6},grid=major,grid style={gray!20}]
\addplot[very thick,mlblue,domain=-6:6,samples=60]{1/(1+exp(-x))};
\addplot[dashed,mlred] coordinates {(-6,0.5)(6,0.5)};
\end{axis}
\end{tikzpicture}
\end{drawit}

\begin{numer}[: logistic regression]
$\theta_0=-4$, $\theta_1=1$, $x=5$ hours. $z=-4+5=1$, $p=\sigma(1)=\frac1{1+e^{-1}}=\mathbf{0.731}\Rightarrow$ pass.\\
Odds $=0.731/0.269=2.72=e^1$. If $y=1$: log-loss $=-\ln0.731=0.313$.
\end{numer}
\begin{numer}[: softmax]
$z=(2,1,0)$: $e^z=(7.389,2.718,1)$, sum $=11.107$ $\Rightarrow$ $P=(0.665,\,0.245,\,0.090)$; predict class 1.
\end{numer}

\apa{outputs probabilities, simple, interpretable (odds ratios), fast}{linear boundary only, needs feature engineering for complex data, sensitive to outliers}{credit default, disease prediction, spam detection, click-through prediction}

\begin{center}
\small
\begin{tabular}{@{}p{3cm}p{6.6cm}p{6.6cm}@{}}
\toprule
 & \textbf{Linear regression} & \textbf{Logistic regression}\\
\midrule
Task & regression (number) & classification (class/probability)\\
Output & $(-\infty,\infty)$ & $(0,1)$\\
Function & identity & sigmoid\\
Cost & squared error (MSE) & log-loss\\
Solution & closed form or GD & GD only\\
\bottomrule
\end{tabular}
\end{center}

\section{Generalised Linear Models (GLMs)}
\begin{dfn}{Generalised linear model}
A \textbf{GLM} extends linear regression to targets that are not normally distributed. It has three parts: (1) a \textbf{random component}: $y$ follows an exponential-family distribution (normal, Bernoulli, Poisson, \dots); (2) a \textbf{linear predictor} $\eta=\vect\theta^{\!\top}\vect x$; (3) a \textbf{link function} $g$ with $g(\mu)=\eta$, where $\mu=\E[y]$.
\end{dfn}
\begin{center}
\small
\begin{tabular}{@{}llll@{}}
\toprule
\textbf{Model} & \textbf{Distribution of $y$} & \textbf{Link $g(\mu)$} & \textbf{Example}\\
\midrule
Linear regression & Normal & identity $\mu$ & price\\
Logistic regression & Bernoulli & logit $\ln\frac\mu{1-\mu}$ & pass/fail\\
Poisson regression & Poisson & log $\ln\mu$ & number of calls per hour\\
\bottomrule
\end{tabular}
\end{center}
\begin{numer}[: Poisson regression]
$\ln\mu=0.5+0.3x$; $x=2$: $\eta=1.1$, $\mu=e^{1.1}=3.0$ calls. $P(y=0)=e^{-3}=0.0498$. Each extra unit of $x$ multiplies the rate by $e^{0.3}=1.35$.
\end{numer}
GLMs are fitted by \textbf{maximum likelihood}.

\section{Support Vector Machines (basics)}
\begin{dfn}{Support vector machine}
A \textbf{support vector machine (SVM)} is a supervised classifier that finds the \textbf{hyperplane} $\vect w\cdot\vect x+b=0$ separating the two classes with the \textbf{maximum margin}. The training points closest to the hyperplane, which define the margin, are the \textbf{support vectors}.
\end{dfn}
\begin{itemize}
  \item Labels $y_i\in\{-1,+1\}$. Margin width $=\frac2{\|\vect w\|}$.
  \item \textbf{Hard margin:} $\min\frac12\|\vect w\|^2$ s.t.\ $y_i(\vect w\cdot\vect x_i+b)\ge1$ (only for separable data; sensitive to outliers).
  \item \textbf{Soft margin:} allow slack $\xi_i$: $\min\frac12\|\vect w\|^2+C\sum\xi_i$. Large $C$ = few violations, narrow margin; small $C$ = wide margin, more violations.
  \item \textbf{Hinge loss:} $\max(0,\,1-y\,f(\vect x))$.
\end{itemize}

\begin{drawit}{Maximum-margin hyperplane}
\centering
\begin{tikzpicture}[scale=0.6]
\draw[-Stealth] (0,0)--(8,0) node[below]{$x_1$}; \draw[-Stealth] (0,0)--(0,6.5) node[left]{$x_2$};
\draw[very thick,mlpurple] (1,6)--(6.5,0.5) node[right,font=\scriptsize]{$\vect w\cdot\vect x+b=0$};
\draw[dashed] (0,5)--(5,0); \draw[dashed] (2,7)--(7.5,1.5);
\foreach \x/\y in {0.8/1.5,1/2.8,2/1.4,2.5/2.5,1.2/0.8}{\fill[mlblue] (\x,\y) circle(4pt);}
\foreach \x/\y in {5.5/3.5,6.5/2.8,5.8/4.8,7/4,4.5/4.5}{\fill[mlorange] (\x,\y) circle(4pt);}
\draw[thick] (2.5,2.5) circle(7pt); \draw[thick] (4.5,4.5) circle(7pt);
\draw[Stealth-Stealth,mlred] (2.5,2.5)--(3.5,3.5) node[midway,above left,font=\scriptsize]{margin};
\node[font=\scriptsize,right] at (8,5.8) {circled = support vectors};
\end{tikzpicture}
\end{drawit}

\begin{numer}[: margin]
Points $(1,1)$ with $y=-1$ and $(3,3)$ with $y=+1$. Solution $\vect w=(0.5,0.5)$, $b=-2$.\\
Check: $(3,3)\to1.5+1.5-2=+1$; $(1,1)\to0.5+0.5-2=-1$ (both on the margin). Margin $=\frac2{\|\vect w\|}=\frac2{0.707}=\mathbf{2.83}=$ distance between the points. New point $(4,1)$: $2.5-2=0.5>0\Rightarrow+1$.
\end{numer}

\section{Binary Classification}
\begin{itemize}
  \item \textbf{Binary classifier:} two classes (positive/negative), e.g.\ ``digit 5 or not''.
  \item Classifier outputs a \textbf{score}; a \textbf{threshold} turns it into a class. Raising the threshold $\Rightarrow$ precision $\uparrow$, recall $\downarrow$ (\textbf{precision--recall trade-off}).
  \item \textbf{ROC curve:} true positive rate (recall) vs.\ false positive rate ($\frac{FP}{FP+TN}$) for all thresholds. \textbf{AUC} $=1$ perfect, $0.5$ random.
  \item \textbf{Multi-class from binary:} One-vs-Rest ($K$ classifiers) or One-vs-One ($K(K-1)/2$ classifiers).
\end{itemize}

\begin{imp}
Most asked: least-squares line + $R^2$ numerical; gradient descent; logistic regression (sigmoid, odds, log-loss); linear vs.\ logistic table; GLM components; SVM margin and support vectors.
\end{imp}

\begin{iq}
\item \twom Define the learning rate. What if it is too large?
\item \twom What is the sigmoid function? Write its range.
\item \twom What are support vectors?
\item \twom What is a link function?
\item \fivem Compare linear and logistic regression.
\item \fivem Explain batch, stochastic and mini-batch gradient descent.
\item \fivem Explain Ridge and Lasso regularisation.
\item \fivem Explain the ROC curve and the precision--recall trade-off.
\item \tenm Fit a least-squares line to given data, predict a value and compute $R^2$.
\item \tenm Explain logistic regression with its cost function and decision boundary.
\item \tenm Explain generalised linear models with examples.
\end{iq}
